\section{Введение}

Современные языковые модели строят распределение следующей языковой единицы по уже наблюденному префиксу. В обычной схеме вычислительная глубина фиксирована: для каждой позиции текста выполняется одно и то же число преобразований. Такая схема проста, но не учитывает неоднородность языковых контекстов. В одних позициях продолжение почти однозначно: после устойчивого словосочетания или синтаксически жесткой конструкции дополнительные вычисления мало меняют распределение. В других позициях полезная информация извлекается только после нескольких внутренних уточнений, поскольку требуется согласовать дальний контекст, смысловую связь и грамматические ограничения.

Эта работа рассматривает число внутренних уточнений как случайную величину, зависящую от текста и состояния модели. Практический вопрос формулируется так: когда следует остановить рекурсивное уточнение и выдать текущее распределение следующей языковой единицы? 

Исходная рекурсивная идея мотивирована малыми рекурсивными моделями \cite{trm}: один и тот же или близкий по параметрам оператор несколько раз применяет уточнение к скрытому состоянию; в указанной работе основные иллюстрации даны на судоку и задачах абстрактного вывода, тогда как в настоящей работе основным объектом исследования является языковое прогнозирование. В качестве дополнительной задачи рассматривалось прогнозирование ответов на логически сложные пазлы ARC-AGI~\cite{arcagi}: для их решения, как правило, требуется большая глубина вычислений, а в малой рекурсивной архитектуре эта глубина создаётся за счёт последовательного уточнения скрытого состояния.

Данную задачу можно разделить на две. Первая задача --- вероятностная: по информации, доступной на текущем внутреннем шаге, оценить распределение оракульного времени остановки. Вторая задача --- решающая: по этому распределению выбрать момент, когда вычисления следует прекратить. Такое разделение устраняет путаницу между статистическим прогнозом и правилом действия. Оно также допускает случай, когда модель считает, что лучший момент уже был в прошлом; эта информация не теряется, а учитывается правилом остановки.

\subsection{Актуальность темы}

Рост вычислительных затрат современных языковых моделей и практика их сжатия, адаптации и раннего выхода делают актуальным статистически корректный учет неопределенности в числе внутренних шагов. Постановка в терминах условного распределения оракульного времени остановки согласуется с классической теорией последовательных решений \cite{wald1947,shiryaev2007} и с современными работами по адаптивной глубине сетей \cite{grave2016,elbayad2020}.

\subsection{Обзор литературы}

Основы последовательного анализа и оптимальной остановки изложены, в том числе, в монографиях и статьях \cite{wald1947,snell1952,robbins1956,shiryaev2007,peskir2006optimal}; динамическое программирование и управление марковскими процессами обсуждаются в~\cite{howard1960dynamic,bertsekas2005dynamic}. Статистическая теория решений задаёт язык для отделения оценки риска от выбора действия~\cite{ferguson1967}; родственные последовательные постановки принятия решений встречаются и в обучении с подкреплением~\cite{sutton2018reinforcement}. Обучение распределений в глубоких моделях и их вероятностная интерпретация систематически описаны в~\cite{goodfellow2016deep,bishop2006}; для последовательностей ключевым остаются рекуррентные кодировочно--декодировочные схемы и механизмы внимания~\cite{hochreiter1997long,mikolov2013efficient,bahdanau2015neural,sutskever2014sequence,vaswani2017attention,dosovitskiy2021image}, а также масштабные предобученные языковые модели~\cite{devlin2019bert,radford2019language,brown2020language}. На практике стохастическая оптимизация параметров почти всегда опирается на адаптивные методы первого порядка~\cite{kingma2015adam}. Экономию вычислений без полного прохода по глубине исследуют идеи раннего выхода и адаптивной глубины~\cite{grave2016,elbayad2020,teerapittayanon2016branchynet,xin2020deebert}; параметрическая адаптация больших представлений облегчается низкоранговым дообучением~\cite{hu2021}. Задачи, близкие по духу к ARC-AGI, обсуждаются в~\cite{chollet2019,arcagi}. Для связи иллюстративного языкового плана экспериментов с открытыми корпусами удобно ориентироваться на описание WikiText-103~\cite{merity2016}. Компактные рекурсивные модели последовательного уточнения представления рассмотрены в~\cite{trm}.

\subsection{Цель работы}

Цель работы --- построить связную вероятностную постановку выбора времени остановки при рекурсивном прогнозировании речи и описать экспериментальную схему, позволяющую сравнивать модели времени остановки по качеству прогноза и вычислительной стоимости.

Для достижения этой цели решаются следующие задачи.
\begin{enumerate}
  \item Описать речь как случайную последовательность на конечном словаре и задать потери предсказания следующей языковой единицы.
  \item Ввести рекурсивный процесс уточнения распределения и связанную с ним фильтрацию внутренней информации.
  \item Определить оракульное время остановки через сумму ошибки и штрафа за вычисления.
  \item Сформулировать задачу предсказания условного распределения оракульного времени остановки.
  \item Выбрать ограниченный набор из шести семейств распределений на \(\{0,\ldots,K\}\), пригодных для описания ранней, поздней и тяжелохвостой остановки.
  \item Задать пять простых правил остановки, строящихся поверх предсказанного распределения \cite{ferguson1967,bishop2006}.
  \item Описать, как переиспользовать малую предобученную языковую модель, чтобы не обучать всю систему с нуля.
  \item Провести эксперименты на ARC-AGI: серию прогонов по семействам распределений, оценку правил остановки на контрольной выборке, сравнение базового обучения и дополнительного обучения только головы оракула.
\end{enumerate}

\subsection{Научная новизна}

Научная новизна прежде всего в том, что в качестве объекта статистики рассматривается именно условное распределение \(\tau^\star\) на \(\{0,\ldots,K\}\): оценивается полный условный закон оракульного момента остановки по внутренней информации текущего шага, а не сводится к одному допустимому числу шагов или одноступенчатой метрике уверенности. На этой основе возможно связать прогноз распределения с решающей процедурой остановки, использующей не только вариант «остановиться сейчас», но и массы на прошлых и будущих шагах; при желании такую связь можно оформить как явную декомпозицию задачи прогноза полного закона \(\tau^\star\) и задачи выбора момента остановки по этому прогнозу. Такая постановка ближе к байесовским и минимаксным схемам принятия решений \cite{ferguson1967}, чем к одноступенчатым эвристикам уверенности.

\subsection{Объект, предмет и метод исследования}

Объект исследования --- рекурсивная языковая модель, которая для одной позиции текста строит несколько последовательных приближений распределения следующей языковой единицы. Предмет исследования --- распределение времени остановки этого внутреннего процесса и правила выбора момента остановки по текущей информации.

Метод исследования является комбинированным. Теоретическая часть задает вероятностную модель и функции риска. Экспериментальная часть сравнивает семейства распределений и правила остановки на языковой задаче, а также на побочной задаче ARC-AGI.
  
\subsection{Научная и практическая значимость}

Научная значимость работы состоит в том, что адаптивные вычисления переводятся в явную статистическую форму: в противоположность распространённым рецептам раннего выхода в компактных рекурсивных моделях по подходу \cite{trm} и в схемах адаптивной глубины для рекуррентных сетей и сетей на основе механизма внимания \cite{grave2016,elbayad2020}, где критерий остановки, как правило, задаётся эвристикой уверенности или фиксированным расписанием без полного вероятностного описания момента остановки, здесь в явном виде рассматривается условное распределение оракульного времени остановки. Это позволяет анализировать не только среднее число шагов, но и форму распределения: наличие ранней массы, длинного хвоста, нескольких режимов сложности и систематического смещения правила остановки.

Практическая значимость связана с уменьшением вычислительной стоимости языкового прогноза. Если часть позиций можно надежно остановить рано, то среднее число внутренних шагов уменьшается без сильного ухудшения качества. Для малых моделей это особенно важно: дополнительная рекурсия может повышать выразительность, но ее цена должна контролироваться. В работе поэтому отдельно рассматривается переиспользование предобученной малой языковой модели и дополнительное обучение только сравнительно небольших добавочных частей.
